% TOP_PROBLEM: {"problemId": "problem.uniqueness-of-the-kpz-fixed-point-among-scaling-covariant-local-fields", "canonicalTitle": "Uniqueness of the KPZ Fixed Point Among Scaling-Covariant Local Fields", "releaseRank": 452, "primaryDomain": "probability_ergodic_dynamics", "primaryDomainLabel": "Probability, ergodic theory & dynamics", "displayStatus": "open", "releaseStatus": "open"}
% !TeX program = lualatex
% ATTEMPT_STATUS: unresolved
\documentclass[11pt]{article}
\usepackage[margin=25mm]{geometry}
\usepackage{fontspec}
\setmainfont{Latin Modern Roman}
\usepackage{amsmath,amssymb,mathtools}
\usepackage{unicode-math}
\setmathfont{Latin Modern Math}
\usepackage{xurl,hyperref}
\hypersetup{colorlinks=true,urlcolor=blue}
\setcounter{secnumdepth}{0}
\setlength{\parindent}{0pt}
\setlength{\parskip}{5pt}
\setlength{\emergencystretch}{3em}
\begin{document}
\section*{\texorpdfstring{Uniqueness of the KPZ Fixed Point Among Scaling-Covariant Local Fields}{Uniqueness of the KPZ Fixed Point Among Scaling-Covariant Local Fields}}\label{452}
\subsection{Definitions and mathematical statement}
Use the state space $\mathrm{UC}$ of upper-semicontinuous profiles $h:{\mathbb{R}}\to{\mathbb{R}}\cup\{-\infty\}$, not identically $-\infty$, with $h(x)\le C(1+|x|)$ for some $C$. The KPZ fixed point $\mathfrak h(t,x;h_0)$ is the profile-valued process defined by Matetski--Quastel--Remenik's Definition 3.12 and transition formula (4.12).

The uniqueness question of Remark 4.12 specifies locality (4.19) and Theorem 4.5(i), (iii), (iv). Locality requires changes to $h_0$ outside neighborhoods of finitely many observation points to change their joint distribution by $o(t)$ as $t\downarrow0$. Scaling means, for $\alpha>0$,
\[
 \alpha\mathfrak h(\alpha^{-3}t,\alpha^{-2}x;
       \alpha^{-1}h_0(\alpha^2\cdot))
 \overset{\rm law}=\mathfrak h(t,x;h_0).
\]
The other two properties are skew time reversibility, interchanging initial and test profiles with sign reversal, and covariance under simultaneous spatial translation. Their full profile-event formulations are retained by the exact theorem reference. The target is uniqueness in law of the nontrivial field satisfying this package, not scaling invariance alone, and not uniqueness of solutions to the KPZ stochastic partial differential equation.
\par\smallskip\textit{Formulation and concept references:} \href{https://arxiv.org/html/1701.00018v3}{[S1]}.

\subsection{Short English statement}
Do the specified locality and symmetry requirements characterize the KPZ fixed point uniquely among nontrivial space-time fields?
\subsection{Research attempt}
\paragraph{A normalization issue in the specified axioms.}
This attempt was prepared on 27 September 2026. The exact references in
[S1] were inspected: Remark 4.12, locality (4.19), and Theorem 4.5(i),
(iii), (iv). The selected package has a time-calibration ambiguity.
If a field $H$ satisfies it, then $H_c(t,x;h)=H(ct,x;h)$ satisfies the
same package for every fixed $c>0$. For the KPZ fixed point these laws
are distinct. A detailed proof is given below, including a nondegeneracy
check. Thus literal uniqueness with fixed numerical units cannot follow
from only the properties enumerated in the catalogue.

This is a formulation diagnostic, not a claim to settle an intended
uniqueness question with an implicit normalization or identification
under changes of units. The original statement is preserved. The present
attempt remains marked unresolved because it does not establish a
characterization after calibration, nor reconstruct all intended
hypotheses behind the source's informal conjectural remark. It would be
misleading to add a time normalization silently and announce a proof of
that changed question.

Recent uniqueness results must also be distinguished by their inputs.
Dauvergne--Zhang [E1] characterize a coupling producing the directed
landscape when the KPZ fixed-point marginals are already prescribed,
together with additional structural conditions. Their primary abstract
does not assert that the four symmetries/locality requirements here
determine those marginals in the first place. No such theorem is imported
into this attempt.

\paragraph{Write the locality and skew events explicitly.}
For initial profiles $h,\widetilde h$ agreeing on neighborhoods of the
finitely many observed points $x_1,\ldots,x_m$, locality compares their
joint distribution functions at fixed thresholds $a_1,\ldots,a_m$:
\[
 \left|\Pr\{H(t,x_i;h)\le a_i\ \forall i\}
       -\Pr\{H(t,x_i;\widetilde h)\le a_i\ \forall i\}\right|=o(t).
\]
The profiles and neighborhoods remain fixed as $t\downarrow0$. This
does not specify the distribution when the nearby initial data are the
same, nor a coefficient for its evolution in time.

The skew reversal property is the equality of full profile events
\begin{equation}\label{a452:skew}
 \Pr\{H(t,\cdot;g)\le-f\}
       =\Pr\{H(t,\cdot;f)\le-g\}.
\end{equation}
The inequalities are imposed at every spatial point in the event, as in
the source's profile formulation. They are not an assertion that the
fields themselves have identical finite-dimensional distributions after
arbitrary interchange of profiles. Keeping that distinction prevents an
incorrect strengthening of the symmetry.

\paragraph{Lemma 452.1: deterministic changes of time preserve the package.}
Let $H$ satisfy the stated scaling, skew reversal, spatial covariance,
and locality. For a constant $c>0$ define
\[
 H_c(t,x;h_0)=H(ct,x;h_0).
\]
Then $H_c$ satisfies all four properties. If $H$ additionally has the
prescribed initial condition, is a time-homogeneous Markov process, or
has a Feller semigroup, each of these extra properties is preserved too.

\emph{Proof of scaling.} The left side of the scaling identity for $H_c$
is
\[
 \alpha H(c\alpha^{-3}t,\alpha^{-2}x;
                  \alpha^{-1}h_0(\alpha^2\cdot)).
\]
Apply the identity for $H$ at time $ct$. The result in law is
$H(ct,x;h_0)=H_c(t,x;h_0)$. The constants $c$ and $\alpha$ commute;
the identity puts no restriction on $c$.

\emph{Proof of the other symmetries.} Substitute $ct$ for $t$ in
\eqref{a452:skew}. Its two sides become exactly the two required events
for $H_c$. Substitution into spatial covariance similarly leaves the
same spatial shift and initial-profile shift. No change to the space or
height coordinates is required.

\emph{Proof of locality.} Denote the original difference of probabilities
at time $u$ by $r(u)$, so $r(u)/u\to0$. The difference for $H_c$ is
$r(ct)$, and
\[
 \frac{r(ct)}t=c\frac{r(ct)}{ct}\longrightarrow0.
\]
For the exponential estimate discussed in [S1], replacing $t$ by $ct$
merely replaces the positive constant in
$\exp(-C\delta^3/t^2)$ by $C/c^2$.

Finally, $ct\downarrow0$ with $t\downarrow0$, so an initial-data limit
is unchanged. If $P_t$ is the transition semigroup, then
$P_t^{(c)}=P_{ct}$ satisfies
$P_t^{(c)}P_s^{(c)}=P_{c(t+s)}=P_{t+s}^{(c)}$.
The Feller mapping and continuity properties are inherited by the same
substitution.\hfill$\square$

In particular, adding just Markovness to the listed package does not
repair the time-calibration issue. Nor does preservation of a fixed
invariant distribution: if $\mu P_t=\mu$ for all $t$, then
$\mu P_{ct}=\mu$ for all $t$. The Brownian statement elsewhere in
Theorem 4.5 concerns the centered profile
$\mathfrak h(t,x;B)-\mathfrak h(t,0;B)$; this centered-profile law is
also insensitive to the speed change. The uncentered Brownian law is
not invariant, since the evolution adds a random global height.

\paragraph{Lemma 452.2: the KPZ laws in this family are distinct.}
Let $Z=\mathfrak h(1,0;0)$ for flat initial data. The one-point law is
a nondegenerate, finite real law; [S1], Section 4.4, identifies the flat
profile with the Airy$_1$ process, and the corresponding GOE
Tracy--Widom distribution is also used in its one-point estimates.
Only the fact that $Z$ is not almost surely zero is needed here.
Scaling at the flat profile gives
\[
 \mathfrak h(t,0;0)\overset{\rm law}=t^{1/3}Z,
 \qquad H_c(1,0;0)\overset{\rm law}=c^{1/3}Z.
\]

For any finite real random variable $Z$ not concentrated at zero, and
any $a>0$ with $a\ne1$, the laws of $Z$ and $aZ$ differ. Indeed, replace
$a$ by its inverse if necessary to assume $a>1$. Equality in law would
give, for every $r>0$ and $n\ge1$,
\[
 \Pr(|Z|>r)=\Pr(|Z|>a^nr).
\]
The right side tends to zero because $Z$ is finite almost surely. Thus
$\Pr(|Z|>r)=0$ for every positive rational $r$, forcing $Z=0$ almost
surely, a contradiction. It follows that $H_c$ and $H_1$ already differ
at the single observation $(t,x,h_0)=(1,0,0)$ whenever $c\ne1$.
\hfill$\square$

This argument uses no fitted variance or approximate quantile. It rules
out the possibility that the time change is merely a different notation
for exactly the same fixed-coordinate law. If the intended uniqueness
is instead modulo choices of time units, then these examples belong to
one allowed equivalence class and the normalized uniqueness problem is
still entirely open within this attempt.

\paragraph{A second diagnostic: the four properties alone do not impose a semigroup.}
Let $Z$ take the values $-1$ and $1$ with equal probability, and define
on the whole stated profile space
\begin{equation}\label{a452:vertical}
 V(t,x;h_0)=h_0(x)+t^{1/3}Z,
 \qquad -\infty+z=-\infty.
\end{equation}
Adding a finite constant preserves upper semicontinuity and the allowed
linear upper bound. The profile is not identically $-\infty$ if the
initial profile was not. It tends back to its initial profile as
$t\downarrow0$.

Scaling is a pathwise identity if the same $Z$ is used on both sides:
$\alpha(\alpha^{-3}t)^{1/3}Z=t^{1/3}Z$. Spatial covariance is immediate.
For any fixed finite observation set, changing the initial profile away
from those points changes none of their values, so locality holds with
zero error. Finally the event in \eqref{a452:skew} becomes
\[
 t^{1/3}Z\le\inf_x[-f(x)-g(x)],
\]
which is symmetric in $f,g$, with the natural extended-real convention
when the bound is $+\infty$. Thus this field satisfies the literal four
conditions as well.

However its proposed transition kernels do not form a time-homogeneous
semigroup. Starting from the zero profile, one step of length two gives
height $2^{1/3}Z$ at the origin, whose probability of being at most zero
is $1/2$. Two independent kernel steps of length one give $Z_1+Z_2$,
whose probability of being at most zero is $3/4$. Therefore
$P_2\ne P_1P_1$. This is not proposed as an alternative KPZ Markov
process. It shows why an axiom package phrased only as one-time
distributional identities does not automatically include temporal
consistency. The speed-change family from Lemma 452.1 survives even
when that consistency is expressly added.

The identity evolution $V(t,\cdot;h_0)=h_0$ is another degenerate
example. Calling the field ``nontrivial'' presumably excludes frozen
evolution, but that phrase alone does not choose among the positive
speeds of a genuinely evolving Markov process. The nondegeneracy test
above avoids relying on a disputed interpretation of triviality.

\paragraph{A natural noise repair also fails in a simple class.}
Could the vertical example be repaired by taking a process $Z_t$ with
independent stationary increments and the same $1/3$ self-similarity?
Assume finite second moments and $Z_0=0$. Such increments imply
$\operatorname{Var}(Z_n)=n\operatorname{Var}(Z_1)$ at positive integers.
Self-similarity would give
$\operatorname{Var}(Z_n)=n^{2/3}\operatorname{Var}(Z_1)$.
At $n=2$ these equalities force the variance to be zero. Likewise
$\mathbb E Z_n=n\mathbb E Z_1=n^{1/3}\mathbb E Z_1$ forces the mean
to be zero. At every fixed $t$, self-similarity then makes $Z_t=0$
almost surely. With a continuous or cadlag version the entire process
is zero, by first intersecting the probability-one events at rational
times and then using path regularity.

This proves failure of this particular finite-variance additive-noise
repair. It is not a nonexistence theorem for KPZ fluctuations: their
profile dynamics do not consist of adding one independent scalar noise
increment uniformly over all space at every step.

\paragraph{Possible calibrations are additional assumptions.}
One can exclude the family $H_c$ by prescribing the flat one-point law
at time one, because Lemma 452.2 then forces $c=1$. Alternatively the
full affine-covariance formula in Theorem 4.5(vi), which is not among
the three selected symmetry items, contains numerical time coefficients.
For a slope $v$, the standard formula has spatial displacement $vt/2$
and vertical correction $v^2t/4$. For $H_c$ these become $cvt/2$ and
$cv^2t/4$, respectively. Fixing those coefficients addresses this speed
freedom in a way the exponent-only scaling law does not.

These are examples of extra data, not axioms silently inserted into the
catalogue. Prescribing one marginal or an affine coefficient would still
not prove that every other initial profile has the correct transition
law. Conversely prescribing the entire transition formula (4.12) would
define the KPZ process directly and would no longer establish a
characterization from the smaller locality/symmetry package.

\paragraph{An explicit sufficient comparison lemma for a normalized route.}
Suppose two Markov semigroups $P_t,Q_t$ act as contractions on a Banach
space of bounded observables, and satisfy the much stronger estimate
\begin{equation}\label{a452:operator}
 \|P_\delta-Q_\delta\|_{\mathrm{op}}=o(\delta)
 \quad(\delta\downarrow0).
\end{equation}
Then $P_t=Q_t$ for every $t\ge0$.

\emph{Proof.} For $\delta=t/m$, the telescoping identity gives
\[
 P_\delta^m-Q_\delta^m
 =\sum_{j=0}^{m-1}P_\delta^{m-1-j}
                  (P_\delta-Q_\delta)Q_\delta^j.
\]
Contraction bounds its operator norm by
$m\|P_\delta-Q_\delta\|=t\,o(\delta)/\delta\to0$.
The semigroup identities identify the left side with $P_t-Q_t$,
independently of $m$. Therefore it is zero.\hfill$\square$

The estimate in the catalogue is not \eqref{a452:operator}. It compares
two initial profiles for the same evolution near fixed observations;
it does not compare two candidate evolutions uniformly over profiles and
observables. Confusing these statements would erase the main uniqueness
problem. A more flexible generator or martingale-problem route would
still need a common, sufficiently determining domain and a uniqueness
theorem for its evolution. Neither is obtained from the four properties
here. In particular time rescaling multiplies any existing generator by
$c$ while preserving those properties.

\paragraph{Finite diagnostics and outcome.}
The companion checks verify exact finite-profile skew events for the
vertical toy, its scaling under rational cube time changes, and the
$1/2$ versus $3/4$ two-step distribution discrepancy. They are tests of
the counterexample mechanisms and conventions, not simulations of the
KPZ fixed point. Distinctness of the genuine time-rescaled KPZ laws
rests on the proved dilation lemma and the attributed nondegenerate
one-point law, not on a discretized numerical distribution.

\textbf{Outcome: UNRESOLVED as an intended normalized characterization;
the literal listed package has a proved calibration defect.}
The four enumerated properties cannot distinguish positive speeds of the
same nontrivial KPZ Markov evolution. A normalized interpretation requires
an explicit additional convention or quotient by changes of units.
After that issue is addressed, this attempt supplies no proof that
locality and the remaining symmetries identify all transition laws.

\subsection{Sources}
\begin{itemize}
\item[S1]K. Matetski, J. Quastel, and D. Remenik, The KPZ fixed point, Remark 4.12, Acta Mathematica 2021.
\url{https://arxiv.org/html/1701.00018v3}.
\textit{Formulation source. Consulted 24 September 2026: Remark 4.12, locality (4.19), and Theorem 4.5(i), (iii), (iv).}
\item[E1]Duncan Dauvergne and Lingfu Zhang,
\emph{Characterization of the directed landscape from the KPZ fixed point},
arXiv:2412.13032 (2024).
\url{https://arxiv.org/abs/2412.13032}.
\textit{Primary abstract and authors' publication link inspected
27 September 2026. Its prescribed KPZ marginals and coupling hypotheses
are distinguished from the package in [S1].}
\item[E2]K. Matetski, J. Quastel, and D. Remenik, same primary text as [S1],
Theorem 4.5, Remarks 4.11--4.12, and Section 4.4, especially (4.24).
\url{https://arxiv.org/html/1701.00018v3}.
\textit{Full symmetry/locality statements and flat-data marginal
identification inspected 27 September 2026. The time-calibration
diagnostic is proved explicitly above.}

\end{itemize}
\end{document}
